\section{Recommendations}

For mobile-oriented front-end development, the results of this study suggest that rendering strategy should be selected by workload and objective rather than by a general rule. In this prototype, SSR was the stronger choice when the goal was to reduce total client-side energy use, because it produced lower median energy in all three scenarios and the clearest support appeared in the static and massive cases. CSR can still be attractive when a lightweight initial navigation response or a shell-first loading pattern is desirable, but that choice should be evaluated against the additional client-side processing required after hydration.

For future implementations, it is advisable to measure energy and browser metrics together rather than relying on speed metrics alone. The corrected benchmark showed that earlier paint-related behaviour did not always align with lower total client-side energy use. In practical terms, teams should validate rendering decisions on representative target devices, use battery-powered measurements when energy is the primary outcome, and interpret outlier-prone data with medians and non-parametric methods.

For follow-up research, the strongest improvements would be methodological rather than conceptual. Repeating the study on multiple devices, adding direct measurement hardware where possible, and extending the benchmark with user interaction scenarios would strengthen generalisability. It would also be useful to measure transfer volume across the full page lifecycle more explicitly, rather than focusing mainly on the initial navigation response.
